\documentclass[a4paper,12pt]{article}
\usepackage[utf8]{inputenc}
\usepackage[T1]{fontenc}
\usepackage[french]{babel}
\usepackage{geometry}
\geometry{left=2.5cm,right=2.5cm,top=2.5cm,bottom=2.5cm}

\usepackage{lmodern}
\usepackage{xcolor}
\definecolor{titleblue}{RGB}{0,51,140}
\definecolor{TITLEBLUE}{RGB}{0,51,140}

\usepackage{hyperref}

\pagestyle{empty}

\begin{document}

\begin{flushleft}
    \textbf{Ismail AISSAOUI} \\
    Ingénieur d'État en Intelligence Artificielle \\
    Chlef, Algérie \\
    ismail.aissaoui.pro@gmail.com \\
    +213 660 70 77 96 \\
    \href{https://ismail-aissaoui.vercel.app}{ismail-aissaoui.vercel.app}
\end{flushleft}

\begin{flushright}
    À l'attention de la \\
    \textbf{Professeure Lilia GZARA} \\
    Laboratoire DISP \\
    INSA Lyon \\
    69621 Villeurbanne, France \\
    \medskip
    Le 29 Avril 2026
\end{flushright}

\bigskip

\noindent \textbf{Objet : Candidature pour un doctorat en IA Explicable (XAI) pour la Décision Industrielle}

\bigskip

Madame la Professeure,

Spécialisé dans la conception de Jumeaux Numériques et les systèmes d'IA pour l'industrie 4.0, je vous adresse ma candidature pour un doctorat sous votre direction au laboratoire DISP. Vos travaux sur l'IA explicable (XAI), la gestion sémantique des connaissances et l'ordonnancement en temps réel correspondent exactement à mes aspirations scientifiques.

Dans mon projet de fin d'études à la Sonatrach, j'ai développé une architecture de \textbf{Jumeau Numérique (Substitut Génératif Mamba-SSM)} pour turbines à gaz qui place l'interprétabilité au centre du diagnostic. En utilisant des modèles de physique informée (\textbf{PINN}) couplés à des méthodes d'explication par saillance de gradient, je me suis attaché à fournir aux ingénieurs opérationnels une justification physique de l'anomalie détectée en temps réel.

Je suis convaincu que l'avenir des systèmes de production repose sur la confiance entre l'humain et l'IA. Je souhaiterais explorer avec vous le développement de modèles de décision capables de s'adapter dynamiquement aux contraintes de production tout en restant "auditables" via des outils d'IA explicable de pointe.

Je vous joins mon CV détaillé. Je serais honoré de pouvoir échanger avec vous sur vos prochaines thématiques de recherche pour l'année 2026.

Je vous prie d'agréer, Madame la Professeure, l'expression de mes salutations les plus distinguées.

\bigskip

\begin{flushright}
    \textbf{Ismail AISSAOUI}
\end{flushright}

\end{document}
